\documentclass[../LuanVan.tex]{subfiles}
\begin{document}

\begin{center}
    \Large{\textbf{Tóm tắt nội dung luận văn}}\\
\end{center}
\vspace{1cm}

Trong hoạt động bán lẻ, việc nhận biết trải nghiệm của khách hàng tại nhiều điểm chạm có ý nghĩa đối với đánh giá chất lượng phục vụ. Các hình thức khảo sát trực tiếp hoặc sau giao dịch có thể làm gián đoạn quá trình mua sắm, có tỷ lệ phản hồi thấp và chịu ảnh hưởng của thiên lệch khảo sát. Nhận dạng biểu cảm khuôn mặt cho phép ghi nhận nhãn biểu cảm từ hình ảnh, nhưng một quan sát riêng lẻ chưa thể mô tả toàn bộ hành trình. Luận văn đề xuất quy trình liên kết kết quả xử lý ảnh với danh tính khách hàng, thời gian và điểm chạm để hình thành chuỗi trạng thái biểu cảm trong từng lần mua sắm. Các chuỗi được sử dụng để thống kê, nhận diện kiểu diễn biến và truy vết về dữ liệu CRM, không thay thế khảo sát trực tiếp hoặc dùng để đánh giá độ hài lòng khách hàng.

Đóng góp về mặt phương pháp là quy trình tổ chức và phân tích dữ liệu biểu cảm có bảo toàn thứ tự giữa các điểm chạm. Ảnh đầu vào được phát hiện khuôn mặt, phân loại biểu cảm, tạo đặc trưng nhận dạng và đối sánh với mẫu đã đăng ký để tạo bản ghi quan sát. Các bản ghi được liên kết theo khách hàng và lần mua sắm, xử lý dữ liệu thiếu, trùng hoặc xung đột thời gian, rồi lựa chọn nhãn đại diện tại mỗi đoạn điểm chạm để tạo chuỗi. Luận văn kết hợp thống kê phân bố và sự thay đổi biểu cảm với Optimal Matching để đo độ khác biệt giữa các chuỗi, PAM để phân cụm và Average Silhouette Width để lựa chọn số cụm. Kết quả gồm chuỗi medoid, quy mô, tỷ lệ, mức độ phù hợp của thành viên và liên kết truy vết về quan sát ban đầu.

Về mặt thực tiễn và kỹ thuật, luận văn xây dựng một hệ thống thử nghiệm tích hợp quản lý CRM, xử lý ảnh và phân tích chuỗi biểu cảm tại nhiều điểm chạm. Để kiểm nghiệm luồng xử lý, luận văn xây dựng một bộ dữ liệu camera thử nghiệm từ các ảnh thuộc tập KDEF, được phát lại theo lịch định trước và gắn với các điểm chạm tương ứng. Hệ thống được kiểm nghiệm từ các thành phần phát hiện khuôn mặt, nhận dạng khách hàng, phân loại biểu cảm và phân cụm chuỗi đến toàn bộ luồng xử lý đầu cuối. Kết quả cho thấy giải pháp có thể tiếp nhận dữ liệu ảnh, liên kết các quan sát theo từng lần mua sắm, hình thành chuỗi trạng thái và cung cấp kết quả phân tích có khả năng truy vết.

\end{document}
